% !TeX program = xelatex
\documentclass[a4paper,10pt]{article}
\usepackage{xeCJK}
\setCJKmainfont{Noto Sans CJK SC}
\setCJKsansfont{Noto Sans CJK SC}
\setCJKmonofont{Noto Sans Mono CJK SC}
\renewcommand{\abstractname}{摘要}
\renewcommand{\today}{\number\year 年\number\month 月\number\day 日}
\usepackage{amsmath,amssymb,booktabs,geometry,array}
\geometry{a4paper,margin=2.4cm}
\newcommand{\code}[1]{\texttt{#1}}

\title{RPN--CG 的 C++ 实现说明与测试记录}
\author{ROPTLIB 项目}
\date{\today}

\begin{document}
\maketitle
\begin{abstract}
本文依据 \code{Solvers/RPNCG.cpp}、\code{Solvers/RPNCG.h} 及相关问题类和测试入口，记录从 Matlab 思路移植到 ROPTLIB 的 RPN--CG 实现：类关系、接口和状态、外层与两个内层求解过程、临时对象的生命周期，以及现有测试能够证明的结果。未取得的跨算法和跨语言数值不以估计值代替。
\end{abstract}

\section{问题、类关系及接口}
测试对象是 Stiefel 流形上的稀疏 PCA：给定 $B\in\mathbb R^{m\times n}$，求
\begin{equation}
 \min_{X\in\mathbb R^{n\times p},\ X^{\mathsf T}X=I_p}
 F(X)=f(X)+g(X)=-\|BX\|_F^2+\lambda\|X\|_1.
\end{equation}
\code{StieSPCA} 继承 \code{Problem}，提供 $f,g$、欧氏梯度 $-2B^{\mathsf T}BX$、欧氏 Hessian 作用 $Y\mapsto-2B^{\mathsf T}BY$、软阈值近端映射 \code{ProxW} 和其广义导数作用 \code{CalJW}。$L=2\|B\|_2^2$ 在问题构造时通过奇异值分解得到。\code{Stiefel} 负责流形随机初值和回缩；测试通过 \code{ChooseParamsSet4()} 选择 polar 型回缩。

\begin{center}
\begin{tabular}{@{}ll@{}}
\toprule
关系 & 职责 \\ \midrule
\code{Solvers $\to$ SolversNSM $\to$ RPNCG} & 通用计时/计数 $\to$ 非光滑求解器接口 $\to$ 本算法 \\
\code{RPNCG $\to$ Problem (StieSPCA)} & 借用目标、梯度、Hessian 作用和近端算子 \\
\code{Problem $\to$ Manifold (Stiefel)} & 借用约束流形和回缩 \\
\bottomrule
\end{tabular}
\end{center}
此处箭头对前三类表示继承，对后两行表示使用/关联，并非所有权转移。\code{RPNCG(prob,initialx)} 调用父类初始化，设置默认参数并复制初始迭代点；析构函数是默认实现。

\subsection{成员函数与参数}
\begin{center}
\small
\begin{tabular}{@{}p{0.27\linewidth}p{0.65\linewidth}@{}}
\toprule
函数 & 实际作用 \\ \midrule
\code{SetDefaultParams()} & 调用 \code{SolversNSM::SetDefaultParams()} 后设定算法参数。 \\
\code{Run()} & 外层迭代、活动集划分、步长调整、回缩和至多三次回溯，记录部分统计。 \\
\code{FindDir()} & 带正则的半光滑 Newton--CG 解乘子方程，返回近端方向 $v$ 并更新 \code{Dinitial}。 \\
\code{ApplyBB()} & 利用 \code{Prob->EucHessianEta} 计算二阶算子作用。 \\
\code{tCG()} & 在活动坐标上作投影 CG 修正 $w$，返回状态码。 \\
\code{UpdateData()}, \code{PrintInfo()}, \code{PrintFinalInfo()} & 前者当前为空；后两者直接调用非光滑求解器父类的打印函数。 \\
\bottomrule
\end{tabular}
\end{center}
文件内辅助函数 \code{ProxAt}、\code{JacobianAt} 将步长 $t$ 转换为 \code{ProxW/CalJW} 所需的标量权重 $1/t$；\code{Aop/AstarOp} 处理约束线性化，\code{GatherVector} 按索引取子向量，\code{BuildBbar} 构造约束投影矩阵，\code{cmidx} 处理列主序索引。这些是文件作用域函数，不是类成员。

\begin{center}
\small
\begin{tabular}{@{}p{0.29\linewidth}p{0.62\linewidth}@{}}
\toprule
成员（默认值） & 用途 \\ \midrule
$\tau=100$, $\gamma=0.01$, $\vartheta=0.01$ & tCG 中非活动坐标的二次惩罚、能量/曲率判断。 \\
$\rho_1=10^{-3}$, $\rho_2=0.5$ & 回溯充分下降阈值和步长缩减比。 \\
$w_1=1.1$, $w_2=0.9$ & 外层近端步长 $t$ 的放大/缩小系数。 \\
$\theta=0.5$, $\kappa=0.1$ & tCG 残差的超线性/线性停止阈值。 \\
\code{max\_innit}$=200$ & tCG 最大迭代数；每步按活动集大小重新设置。 \\
\code{max\_inner\_iter} & 在头文件声明，但当前实现不赋值、不使用。 \\
继承的 \code{Tolerance}, \code{Max\_Iteration}, \code{Verbose} 等 & 停止参数和输出配置；$x_1,x_2,f_1,f_2$、\code{iter, ComTime, nf, ng, nR, nH} 等由基类存放。 \\
\bottomrule
\end{tabular}
\end{center}
头文件还有独立的 \code{RPNCGParams} 结构体（其中不含 $\tau$）；当前 \code{Run()} 读取的是 \code{RPNCG} 自身的成员，\emph{没有}使用该结构体。\code{RPNCGTCGStatus} 中 $-1$ 表示达到内迭代上限，$0,1,2,3$ 分别是两种初期退出、负曲率和能量判据退出，$4,5$ 分别对应线性和超线性残差判据。

\section{运行逻辑}
\subsection{外层迭代}
\code{Run()} 先调用父类 \code{Run()} 初始化计时/输出缓存，计算 $F(x_1)$ 和欧氏梯度，设 $t=1/L$（仅当问题可转换为 \code{StieSPCA}；否则取 $1$），初始化零乘子矩阵 $D\in\mathbb R^{p\times p}$。在 $\mathrm{err}>\code{Tolerance}$ 且 $\code{iter}<\code{Max\_Iteration}$ 时执行：
\begin{enumerate}
\item 根据前一轮 $\|v\|_F$ 和 tCG 状态收紧 \code{innertol}；调用 \code{FindDir} 求 $v$，以 $\|v\|_F$ 记作 $n_v$。乘子 $D$ 跨外层迭代热启动。
\item 逐个列主序位置 $k$ 判断：若 $(x+v)_k\neq0$ 且 $|x_k|\geq n_v$，放入 \code{id1}，否则放入 \code{id2}；将 tCG 上限设置为 $\max(1,\min(200,\operatorname{round}(1.2|\code{id1}|)))$。
\item 计算 $B_\lambda=xD$、\code{tmp11}$=x^{\mathsf T}B_\lambda$，调用 \code{tCG} 求仅定义在 \code{id1} 上的修正 $w$；拼出 $u_{\code{id1}}=v_{\code{id1}}+w$、$u_{\code{id2}}=v_{\code{id2}}$，令 $n_u=\|u\|_F$。
\item 若 $(4+1/t)n_u<n_v$ 或状态为 $0$，令 $t=\max(t_0,w_2t)$；否则状态 $\leq4$ 时令 $t=w_1t$。状态 $\geq5$ 使用带 \code{flag} 的特殊分支（可能回到上一轮保存的 $x_{\rm old},u_{\rm old}$），普通分支用 $x_2=R_{x_1}(\alpha u)$，$\alpha$ 从 $1$ 开始按 $\rho_2$ 缩减，至多追加三次回溯。
\item 计算新梯度，设 $\mathrm{err}=\min(n_v,n_u)$，更新 $x_1,f_1$ 和梯度，累计 \code{iter}、目标/梯度/回缩调用次数；遇到新目标值 NaN/Inf 则中止。循环结束记录 \code{ComTime}。
\end{enumerate}
这里的停止条件是 \code{Run()} 自行实现的绝对方向范数判据，不调用父类 \code{IsStopped()}；因此基类 \code{Stop\_Criterion}、\code{TimeBound} 和 \code{StopPtr} 不控制此循环。

\subsection{乘子半光滑 Newton--CG：\code{FindDir}}
令 $\operatorname{sym}(Z)=(Z+Z^{\mathsf T})/2$，源码定义
\begin{align}
 A_x(\eta)&=-\tfrac12(\eta^{\mathsf T}x+x^{\mathsf T}\eta),
 & A_x^*(\Lambda)&=-x\operatorname{sym}(\Lambda),\\
 B(D)&=x-t\bigl(\nabla f(x)-A_x^*(D)\bigr),
 & F(D)&=A_x\bigl(\operatorname{prox}_{tg}(B(D))-x\bigr).
\end{align}
以传入的 $D$ 为初值，使用 \code{CalJW} 作用形成 Jacobian--向量积；每个 Newton 步用最多 30 次 CG 求正则化线性方程（正则系数随残差变化），最多 1000 次 Newton 更新。接受条件使用长度为 5 的非单调残差窗口；拒绝时依据 $\rho$ 调整正则系数，并尝试投影修正/残差步。循环判据是 $\|F(D)\|_F^2>\code{innertol}$，返回 $v=\operatorname{prox}_{tg}(B(D))-x$、迭代次数及新乘子；这不是直接求 tCG 的修正量。

\subsection{约束投影 tCG 与二阶算子}
\code{ApplyBB} 对 $Y$ 计算
\begin{equation}
 \mathcal B_x(Y)=\nabla^2 f(x)[Y]+Y(x^{\mathsf T}B_\lambda)
 +x\operatorname{sym}(Y^{\mathsf T}B_\lambda).
\end{equation}
\code{BuildBbar} 对称矩阵空间的 $p(p+1)/2$ 个标准正交基 $S_j$ 构造
$\bar B_{kj}=(xS_j)_{\code{id1}[k]}$；对角基系数为 $1$，非对角基系数为 $1/\sqrt 2$。令 $D_B=\bar B^{\mathsf T}\bar B$，通过伪逆构造投影 $P(z)=z-\bar B D_B^\dagger\bar B^{\mathsf T}z$。tCG 先检查基于 $\mathcal B_x(v)$、$g(x+v)-g(x)$ 的能量条件；若 \code{id1} 为空立即返回空修正。其余情况下在投影子空间从 $w=0$ 开始做 CG，检查负曲率、能量下降及 $\|r_j\|\leq\|r_0\|\min(\|r_0\|^\theta,\kappa)$；未满足而迭代耗尽则返回零修正。注意源码在循环尾更新 \code{inner\_it}，某些提前退出的日志计数可能小于实际完成的 CG 步数，不能据此精确核算 Hessian 次数。

\section{内存管理与统计边界}
\code{Prob}、\code{Mani} 是基类保存的非拥有型指针：调用方须确保问题和流形比求解器活得更久；测试中三者均为栈对象。\code{x1/x2}、乘子 $D$、子向量和中间矩阵使用 \code{Variable/Vector} 值对象及局部 \code{std::vector<integer>} 管理，函数返回值和容器按 C++ 生命周期释放，不在 \code{RPNCG} 中手动 \code{new/delete}。\code{BuildBbar} 直接生成大小为 $|\code{id1}|\times p(p+1)/2$ 的矩阵，避免显式存储 Matlab 风格的三维中间张量；另需保存 $p(p+1)/2$ 阶的 Gram 矩阵及其伪逆。\code{GatherVector} 会复制活动集数据，\code{ApplyBB} 及矩阵乘法也会产生临时 \code{Vector}，不能声称整体零拷贝。

基类的 \code{nf}、\code{ng}、\code{nR} 分别在 \code{Run()} 的目标值、梯度、回缩调用点显式自增；但 \code{FindDir} 内 \code{ProxW/CalJW}、tCG 内 \code{g()} 以及 \code{ApplyBB} 中 \code{EucHessianEta} 的调用未通过专用计数器完整统计。尤其 \code{nH} 没有在该实现中增加，不能把打印的零当作没有二阶计算。

\section{测试记录与对比边界}
\subsection{现有 C++ 用例}
入口 \code{test/mainRPNCG.cpp} 使用随机种子 $2$；\code{test/testRPNCG.cpp} 取 $(n,m,p)=(800,50,5)$、$\lambda=0.8$，对高斯矩阵每列去均值并归一化，再随机生成 Stiefel 初值。设置 $\code{Max\_Iteration}=1000$、$\code{Tolerance}=10^{-6}$、\code{Verbose=ITERRESULT}，运行一次 \code{solver.Run()}。对仓库已有 \code{build/bin/testRPNCG} 可执行文件的一次实际运行记录如下（程序逐步打印计时，因此仅作单次样例，非稳定基准）：
\begin{center}
\begin{tabular}{@{}lrrrrr@{}}
\toprule
方法 & 外层步数 & 最后输出的 $F$ & 程序内用时 (s) & 最后 $n_v$ & 最后 $n_u$ \\ \midrule
RPN--CG (已有 C++ 二进制) & 147 & $-3.850512\times10^{1}$ & $0.558956$ & $8.693090\times10^{-8}$ & $1.513399\times10^{-6}$ \\
\bottomrule
\end{tabular}
\end{center}
末次 $\min(n_v,n_u)<10^{-6}$，故按源码循环判据退出；这是\emph{现有二进制}的观察值，未验证其编译内容是否与当前源码完全一致。日志未给出 \code{nf/ng/nR} 的最终总数，也没有可靠的 Hessian/近端调用次数，故这些数值暂记为\textbf{未记录}，而非由迭代次数推断。

\subsection{同口径比较所需补测}
\begin{center}
\small
\begin{tabular}{@{}p{0.31\linewidth}p{0.15\linewidth}p{0.19\linewidth}p{0.25\linewidth}@{}}
\toprule
比较对象 & 时间 (s) & 主要计算次数 & 效果/可比性 \\ \midrule
C++ RPN--CG（上述用例） & $0.558956$（单次） & 外层 $147$；\code{nf/ng/nR} 未记录 & $F\approx-38.50512$ \\
C++ 一阶方法（如 \code{IARPG/IRPG}） & 未取得 & 未取得 & 未在同一数据和初值下运行 \\
Matlab 二阶 RPN--CG & 未取得 & 未取得 & 仓库中未找到对应的 Matlab RPN--CG 驱动程序/结果 \\
\bottomrule
\end{tabular}
\end{center}
仓库 \code{test/TestStieSPCA.cpp} 中虽有 \code{IARPG} 测试，但默认 $(n,m,p)=(1000,20,5)$、$\lambda=2$，与本用例不同，不能直接做效率比；\code{Matlab/ForCpp/CTestStieSPCA.m} 配置的是一阶方法（$n=1000,p=17$ 等），不是所要求的 Matlab 二阶基线。当前会话尝试启动 Matlab 获取版本信息时未在 8 秒内返回，因此没有可验证的 Matlab 二阶运行值。

为完成 \code{req.txt} 的两项对比，应固定\emph{同一个} $B$ 和 $X_0$（将输入数据及初值在两端共享），统一 $\lambda$、回缩、精度、最大迭代数和停止判据，并在相同硬件/线程数下分别多次计时（报告中位数）。至少记录最终目标值、$\min(n_v,n_u)$ 或共同定义的残差、外层次数、总耗时、目标/梯度/回缩次数，以及 Hessian--向量积、近端映射、Newton 和 CG 的实际总调用次数。对 Matlab 需先取得相应二阶代码及可运行环境；对 C++ 需先在同一测试入口运行一阶求解器。若仅允许修改 TeX，这些测试入口和计数埋点不能在本次工作中补写；故不作无依据的速度优劣结论。

\end{document}
